% Part: intuitionistic-logic
% Chapter: tableaux
% Section: introduction

\documentclass[../../../include/open-logic-section]{subfiles}

\begin{document}

\olfileid{int}{tab}{int}

\olsection{مقدمة}

!!^{tableau} أشجار مخصوصة تتفرع إلى أسفل، وعقدها من !!{signed
  formula}؛ وهي أزواج يجتمع في كل واحد منها علامة لقيمة الصدق
($\True$ أو~$\False$) و!!a{formula}:
\[
\sFmla{\True}{!A} \text{ أو } \sFmla{\False}{!A}.
\]
ويُفتتح !!^a{tableau} بجملة من \emph{الافتراضات}، ثم لا ترد فيه
!!{signed formula} بعدها إلا عن تطبيق قاعدة من قواعد الاستدلال. فمن
القواعد ما يزيد طرفًا قائمًا بواحدة أو أكثر من !!{signed formula}،
ومنها ما يُنشئ طرفين فيتشعب منه الفرع. وما تنتجه القاعدة من
!!{signed formula} تكون !!{formula} فيه أدنى تركيبًا من نظيرتها في
!!{signed formula} التي وقعت عليها القاعدة. فإذا اجتمع في شعبة
$\sFmla{\True}{!A}$ و$\sFmla{\False}{!A}$ سميناها \emph{مغلقة}؛
فإن أُغلقت شعب !!a{tableau} كلها سُمّي !!{tableau} نفسه مغلقًا.
والـ!!{tableau} المغلق !!a{derivation} يشهد بأن مجموعة
!!{signed formula} التي بُدئ بها !!{tableau} ممتنعة الإرضاء. ومن هذا
نحدّ العلاقة~$\Proves$: فنقول
${\OLId{greek-variable}{Gamma}} \Proves !A$ إذا وفقط إذا وجدت مجموعة منتهية
${\OLId{greek-variable}{Gamma}}_{\OLMathNumeral{0}} = \{!B_{\OLMathNumeral{1}}, \dots, !B_{\OLId{latin-ordinary}{n}}\} \subseteq {\OLId{greek-variable}{Gamma}}$ بحيث يوجد
!!{tableau} مغلق للافتراضات
\[
\{\sFmla{\False}{!A}, \sFmla{\True}{!B_{\OLMathNumeral{1}}}, \dots, \sFmla{\True}{!B_{\OLId{latin-ordinary}{n}}}\}.
\]

وأما في المنطق الحدسي فلا بد من توسيع معنى !!{signed formula} وتغيير
قواعد الروابط تبعًا لذلك. إذ تحمل !!{formula} في !!{tableau} الحدسية،
مع العلامة ($\True$ أو~$\False$)، \emph{سابقة}~$\sigma$ أيضًا.
والسابقة متتالية غير خالية من الأعداد الصحيحة الموجبة، أي
$\sigma \in (\PosInt)^* \setminus \{\emptyseq\}$. نكتب هذه السوابق
من غير~$\tuple{\ }$ تحيط بها، ونجعل بين !!{element} المفردة نقاطًا~$.$
لا فواصل~$,$. فإذا كانت~$\sigma$ سابقة، دلّت~$\sigma.{\OLId{latin-ordinary}{n}}$ على
$\sigma \concat \tuple{{\OLId{latin-ordinary}{n}}}$؛ فإذا كانت مثلًا
$\sigma = {\OLMathNumeral{1}}.{\OLMathNumeral{2}}.{\OLMathNumeral{1}}$، كانت~$\sigma.{\OLMathNumeral{3}}$ هي~${\OLMathNumeral{1}}.{\OLMathNumeral{2}}.{\OLMathNumeral{1}}.{\OLMathNumeral{3}}$. ومن ثم، على سبيل المثال،
\[
\sFmla{\True}{!A \lif (!B \lif !C)}[{\OLMathNumeral{1}}.{\OLMathNumeral{2}}]
\]
تُسمّى \emph{!!{signed formula} ذات سابقة} (ويقال اختصارًا
\emph{!!{formula} ذات سابقة}).

وعلى التأويل الحدسي تشير السابقة إلى عالم من نموذج يجوز أن يرضي
!!{formula} الواقعة في شعبة من !!a{tableau}. فإذا أشارت~$\sigma$
إلى عالم، أشارت~$\sigma.{\OLId{latin-ordinary}{n}}$ إلى عالم يمكن بلوغه من العالم الذي تشير
إليه~$\sigma$.

وعلاقة الإتاحة في النماذج الحدسية انعكاسية متعدية. ويقابل ذلك في لسان
السوابق أن تكون~$\sigma$ متاحة من ذاتها، وأن تكون كل سابقة تمد~$\sigma$
متاحة كذلك،
أي كل سابقة صورتها~$\sigma.{\OLId{latin-ordinary}{n}}_{\OLMathNumeral{1}}.\cdots.{\OLId{latin-ordinary}{n}}_{\OLId{latin-ordinary}{k}}$. ونضع الرمز~$\sigma.*$
للدلالة على~$\sigma$ وعلى أي امتداد منها؛ فالسوابق التي يدل عليها
$\sigma.*$ هي السوابق المتاحة من~$\sigma$ كلها، ولا يدخل فيها غيرها.

\end{document}


